\chapter{Validation, source map and reproducibility}
\label{appendix:validation}

\section{What was checked for this edition}
The implementation reference and source links are pinned to revision
\texttt{\SourceShort}; its full identifier is recorded in \path{edition.tex}.
The original tutorial and recovery examples were
checked with newly built Linux x64 programs from the earlier baseline
\texttt{9bedab0c9de307d0435d3abfb3eab98a2f6d42a5}, using
MiniLangCompilerPy 1.3.10 at compiler revision
\texttt{8e64c9cfda92199b68ac5d3974fea445f12f6c39}.
The server, native client, consistency checker and backup tool were compiled
for the validation. Each run used fresh, disposable database directories.
No existing user database was used as a mutable test fixture.

The manual validation runner executes four complete SQL files through the
native console and checks 53 focused outcomes through the repository's Python
connector. A separate recovery runner verifies the offline backup/restore
procedure, runs the native checker before backup and after restore, starts
the restored database, compares counts and a joined aggregate, and verifies
a writable transaction followed by rollback.

The initial stored-procedure language expansion retains its separate evidence
at revision \texttt{139eb3c}. The subsequent five-package expansion at this
edition's implementation passes 12 focused native Linux cases, including
127 runtime, 21 edge, 154 security/metadata and 105 final parity assertions.
The complete focused run included 100 parity assertions; a later targeted
105-assertion run verifies the final character-set/collation SQLSTATE correction.
Native protocol and authenticated-protocol regressions also pass. Six routine
test programs cross-compile to PE32+ AMD64, and the final parity program is
rebuilt separately after that correction. Windows server, console and
Workbench entry points also cross-compile with validated PE32+ AMD64 images.
No Windows executable was run.

Connector checks pass 18 Python unit tests, 16 general JDBC checks and
77 procedure-specific JDBC checks. Java execution uses Java 21 with
source/target 11 compilation; it is not Java 11 runtime acceptance. A freshly
rebuilt native server passes 33 top-level connector checks, one of which runs
the 25-check JDBC live suite. These two live counts are not additive.
The complete incremental evidence, source hashes and final application-build
results are in \path{docs/acceptance/stored-procedures-2026-11-25.json}.
Historical tutorial and initial procedure records have not been relabeled as
new executions at the current revision.

\begin{longtable}{p{.29\linewidth}p{.63\linewidth}}
\caption{Documentation validation and its limits}\\
\toprule Area & Evidence and interpretation\\\midrule\endhead
Teaching setup & Four tables and 19 data rows loaded through the native script client. Counts are 4 customers, 5 products, 4 orders and 6 order items.\\
Queries & Filtering, sorting, NULL counts, joins, grouped totals, outer-join anti-matching, CTEs, recursion, windows, working subquery shapes and expected results.\\
Changes & Insert/update/delete with RETURNING, savepoint rollback, full rollback, duplicate-key rejection, missing-parent rejection and CHECK rejection.\\
Known limitations & Decimal scaling/presentation, temporal comparison, CURRENT\_TIMESTAMP, NULLIF, legacy text NULL ambiguity, unsupported DELETE CASCADE, orphan-producing parent-key updates, correlated NOT EXISTS, identity reuse and per-row sequence materialization.\\
Operational recipe & Clean stop, native checker, offline backup, restore to a new path, checker on the restored copy, restarted SQL and reversible write.\\
Source review & Parser/binder/executor, command entrypoints, config, security, drivers, Workbench, planner, storage, recovery and replication implementations.\\
Not a claim of execution & This work did not run the full Windows acceptance suite, Windows Workbench, all TLS certificate scenarios, long-duration soak, every replication failure scenario or every SQL combination. Those descriptions are source-reviewed and linked to relevant tests.\\\bottomrule
\end{longtable}

\warning{A passing limitation probe is not a correctness endorsement}{Some
checks deliberately assert an observed defect or unsupported behavior. Such a
check passes when the manual describes the running implementation faithfully.
For example, confirming that \code{NULLIF(5,NULL)} returns null does not make
that result standard-conforming. When the engine changes, a failed limitation
probe may indicate that a warning can be removed after investigation.}

\section{Reproduce the teaching checks}
From the repository root, with matching native executables already built:
\begin{lstlisting}[language=bash]
python3 docs/manual/tools/validate_examples.py \
  --bin-dir build/bin-linux \
  --output build/manual-validation.json

python3 docs/manual/tools/validate_recovery.py \
  --bin-dir build/bin-linux \
  --output build/manual-recovery-validation.json
\end{lstlisting}
On Windows use \code{python} and \path{build/bin}. The runners discover a
free loopback port, create temporary directories under \path{build/}, start
and stop their own processes, and remove their disposable databases when
finished. They do not need a running application server and must not be
pointed at a production database. \code{--bin-dir} selects executables, not
a database path.

The committed \path{validation.json} and \path{recovery-validation.json}
retain the original tutorial/recovery outcomes. They record their source baseline,
platform, script output, expected values and observed results. Rerun with an
output under \path{build/} while comparing a new revision; replace the
committed evidence only as part of a reviewed documentation update.

\section{Selected observed edge cases}
The table below preserves the tutorial-baseline probes. Read the corresponding
language and operations chapters for current explanations and practical
workarounds. In the present procedure expansion, negotiated results-v2 fixes
the literal-NULL ambiguity, and cursor admission no longer evaluates NEXTVAL
once during a probe and again during fallback. Those new regression results
are recorded separately; the historical observations below are unchanged.

\begin{longtable}{p{.45\linewidth}p{.47\linewidth}}
\caption{Recorded tutorial-baseline boundary cases reviewed for this edition}\\
\toprule Probe & Observed outcome\\\midrule\endhead
\code{CAST('3.30' AS DECIMAL(10,2))} & Wire cell \code{330}; the schema scale is not applied by the generic formatter.\\
Add decimal 3.30 and 1.20 & Wire cell \code{45000}; arithmetic adds raw coefficients and rescales incorrectly.\\
Compare decimal 3.30 to integer 330 & \code{TRUE}; mixed-scale comparison is not a reliable numeric comparison.\\
\code{NULLIF(5,NULL)} & SQL NULL.\\
Compare two \code{CAST(0 AS DATE)} values & Error: values are not comparable.\\
\code{CURRENT\_TIMESTAMP} & Wire cell \code{0}, not the current wall-clock time.\\
\code{5 / 2} & \code{2.5}; division deliberately produces a floating result.\\
\code{SELECT 'NULL', NULL} through Python & Both values appear as Python \code{None}.\\
Delete referenced parent with DELETE CASCADE & Rejected with an implementation-stage diagnostic.\\
Update parent key from 1 to 2 with existing child & Update succeeds; child remains at 1.\\
Correlated EXISTS and scalar COUNT subquery & Correct expected fixture results in the tested shapes.\\
Correlated NOT EXISTS & Runtime error 1306: \code{Struct has no member 'kind'}.\\
Sequence starts at 1; query NEXTVAL over products ordered by key & Five rows all display 2; evaluation is not once per row and the cursor/fallback path consumes twice.\\
Next standalone NEXTVAL on that sequence & 4, reflecting the same double-materialization path.\\
BEGIN READ ONLY; NEXTVAL; ROLLBACK & NEXTVAL succeeds with 6; CURRVAL remains 6 after rollback.\\
Identity rows 1,2; delete 2; insert again & The new row reuses 2 because allocation uses the current maximum.\\\bottomrule
\end{longtable}

\section{How the manual covers the repository}
\begin{longtable}{p{.28\linewidth}p{.64\linewidth}}
\caption{The implementation areas behind the handbook}\\
\toprule Area & Primary files and corresponding chapters\\\midrule\endhead
Build and applications & \path{build.ps1}; all six \path{src/apps/*/main.ml} entrypoints. Installation, console, administration and recovery chapters.\\
SQL surface & \path{src/minisql/sql/parser.ml}, \path{lexer.ml}, \path{ast.ml}, \path{binder.ml}, \path{types.ml}, \path{values.ml}, \path{expressions.ml}. Language chapters.\\
Execution & \path{src/minisql/executor/executor.ml}, \path{dml.ml}, aggregate, join, projection and sort modules. Language, transaction and performance chapters.\\
Catalog and schema & \path{src/minisql/catalog/catalog.ml}, \path{schema_history.ml}, \path{statistics.ml}. Definitions, permissions, metadata and internals.\\
Transactions & \path{src/minisql/transaction/transaction.ml}, WAL, recovery and checkpoint modules; database manager and logical locks. Transaction and recovery chapters.\\
Server and configuration & \path{src/minisql/server/}, \path{src/minisql/config/}, logger, diagnostics and deployment templates. Installation and administration.\\
Security & Catalog security, common crypto helpers, TLS platform adapters, protocol authentication, TDE and backup encryption. Security and integration.\\
Client interfaces & \path{src/minisql/client/}, \path{src/minisql/admin/}, \path{clients/python/}, \path{clients/jdbc/}. Clients chapter.\\
Optimizer & \path{src/minisql/planner/}, statistics and executor operators. Performance and queries.\\
Physical formats & \path{src/minisql/storage/}, \path{docs/formats/}, \path{common/format_version.ml}. Storage chapter.\\
Operational tools & \path{src/minisql/tools/}, \path{tools/monitoring/}, \path{tools/replication/}, recovery/fault/soak tests. Administration, recovery and diagnostics.\\\bottomrule
\end{longtable}

The abbreviated module names in this overview inherit the directory of their
first fully qualified sibling. Exact links in the chapters always identify
the full repository path. \path{source-map.json} lists the verified source
references by chapter so that an update can identify which chapters depend
on a changed file. The generated API documentation in \path{docs/api/}
complements this book with declaration-level detail.

\section{Build the PDF from its LaTeX sources}
Install a normal TeX Live or MiKTeX distribution providing pdfLaTeX,
MakeIndex, Latin Modern, TikZ/PGF, tcolorbox, listings, booktabs, longtable,
tabularx, enumitem, titlesec, fancyhdr, microtype, hyperref and bookmark.
The commonly supplied recommended and extra LaTeX packages contain these.
The build uses no shell escape and needs no raster-image generator.

\begin{lstlisting}[language=bash]
python3 docs/manual/tools/build_manual.py
\end{lstlisting}
The build writes auxiliary files to \path{build/manual/}, runs the index
and enough LaTeX passes to resolve navigation, and copies the final document
to \path{docs/manual/minisql-user-manual.pdf}. To choose another output path,
use \code{--output}. All illustrations are editable TikZ source files under
\path{figures/}; all tutorial SQL is under \path{examples/}. The main document
is \path{minisql-user-manual.tex}, with formatting in \path{preamble.tex} and
the pinned source identity in \path{edition.tex}.

\section{Review checklist for the next edition}
When the implementation changes, inspect the changed code and identify
affected manual chapters through the source map. Rebuild native programs
with a recorded compiler. Run the tutorial and recovery checks, then revise
the SQL syntax, limits, examples and known-defect descriptions as needed.
Keep historical release evidence clearly separate from new validation.
Rebuild the PDF and inspect pages containing long command lines, dense
tables, diagrams and chapter transitions. Check hyperlinks, contents entries,
the index, and the absence of missing references.

Publish the source, PDF, executable examples and validation evidence together.
A manual is most useful when another person can both read its explanations
and reproduce the behaviors it describes.
\sourcecode{docs/acceptance/STORAGE_V2_2026-11-24.md}{The separate storage-v2 implementation gate and its explicit platform scope.}
\sourcecode{tests/run_storage_v2.py}{Native focused and selected storage regression runner.}
